\section{Proofs}
\label{app:proofs}

This appendix collects the proofs deferred from
Section~\ref{sec:method}. We use throughout the standard logistic
identities $\sigma(t):=(1+e^{-t})^{-1}$, $\softplus(t):=\log(1+e^{t})$,
with $\softplus'(t)=\sigma(t)$, $\softplus(t)-\softplus(-t)=t$,
$\sigma(-t)=1-\sigma(t)$, and $\sigma(t)(1-\sigma(t))=\softplus''(t)>0$.
The binary entropy in nats is
$H_2(p):=-p\log p-(1-p)\log(1-p)$ for $p\in(0,1)$. All claims below are
stated for the proxy abstention loss $L$ in~\eqref{eq:proxy-target},
the surrogate $\Phi_{\mathrm{abs}}$ in~\eqref{eq:surrogate}, and the
pointwise minimiser $s^\star$ of Proposition~\ref{prop:pointwise-min}.

\subsection{Proof of Lemma~\ref{lem:bayes-rej}}
\label{app:bayes}

Fix $x\in\Xcal$ and condition on $X=x$. By Tonelli's theorem the
conditional risk of the abstention loss \eqref{eq:abs-loss} at $x$
decomposes as
\[
  \E\!\bigl[\ell_{\mathrm{abs}}(h,r;X,Y)\bigm|X=x\bigr]
  \;=\;\eta_h(x)\,\ind\{r(x)=0\}\;+\;c(x)\,\ind\{r(x)=1\},
\]
since $h$ and $r$ are deterministic given $x$. The pointwise infimum
over $r(x)\in\{0,1\}$ is $\min\{\eta_h(x),c(x)\}$, attained by $r(x)=0$
when $\eta_h(x)\leq c(x)$ and by $r(x)=1$ when $\eta_h(x)>c(x)$
(either choice is optimal on the measure-zero set
$\{x:\eta_h(x)=c(x)\}$). Aggregating in $x$ and applying the tower
property gives the Bayes risk
$\E_X[\min\{\eta_h(X),c(X)\}]$, attained by every measurable $r^\star$
that satisfies the stated pointwise rule. \qed

\subsection{Proof of Proposition~\ref{prop:pointwise-min}}
\label{app:pointwise-min}

Fix $x\in\Omega$ with $\rho^\theta(x),c(x)>0$ and write
$\rho\equiv\rho^\theta(x)$, $c\equiv c(x)$. Differentiating
\eqref{eq:surrogate} in $s$,
\[
  \partial_s\Phi_{\mathrm{abs}}(s;x)
  \;=\;-\rho\,\sigma(-s)+c\,\sigma(s)
  \;=\;(\rho+c)\,\sigma(s)-\rho,
\]
where the second equality uses $\sigma(-s)=1-\sigma(s)$. Setting the
derivative to zero gives $\sigma(s^\star)=\rho/(\rho+c)$, hence
$s^\star=\log(\rho/c)$. The second derivative,
$(\rho+c)\,\sigma(s)(1-\sigma(s))$, is strictly positive because
$\rho+c>0$ and $\sigma(s)(1-\sigma(s))>0$ for all $s\in\R$, so
$\Phi_{\mathrm{abs}}(\,\cdot\,;x)$ is strictly convex. Coercivity follows
from $\softplus(t)\sim t$ as $t\to+\infty$ and $\softplus(-t)\sim -t$ as
$t\to-\infty$, so $\Phi_{\mathrm{abs}}(s;x)\to+\infty$ at both ends and
$s^\star$ is the unique global minimiser. The sign identity
$s^\star(x)>0\Leftrightarrow\rho^\theta(x)>c(x)$ is immediate from
$s^\star=\log(\rho/c)$. \qed

\subsection{Proof of Theorem~\ref{thm:h-consistency}}
\label{app:hconsistency}

We first establish a pointwise lower bound on the surrogate excess
$\delta(x):=\Phi_{\mathrm{abs}}(s(x);x)-\Phi_{\mathrm{abs}}(s^\star(x);x)$
on the disagreement set $S:=\{x:r_s(x)\neq r_{s^\star}(x)\}$, then
combine it with a Cauchy--Schwarz step to obtain
\eqref{eq:h-consistency-main}. Throughout this subsection we abbreviate
$\rho\equiv\rho^\theta(x)$ and $c\equiv c(x)$ at a generic point and
recall that $\rho,c>0$ on the active grid by the standing assumption of
Section~\ref{sec:method-surrogate}.

\paragraph{Value of the surrogate at $0$ and at $s^\star$.}
By \eqref{eq:surrogate} and $\softplus(0)=\log 2$,
\begin{equation}
  \Phi_{\mathrm{abs}}(0;x)\;=\;\rho\,\log 2+c\,\log 2\;=\;(\rho+c)\,\log 2.
  \label{eq:phi-zero}
\end{equation}
Substituting $s^\star=\log(\rho/c)$ into \eqref{eq:surrogate} and using
$\softplus(\log(\rho/c))=\log(1+\rho/c)=\log\!\frac{\rho+c}{c}$ and
$\softplus(-\log(\rho/c))=\log\!\frac{\rho+c}{\rho}$,
\begin{equation}
  \Phi_{\mathrm{abs}}(s^\star;x)
  \;=\;\rho\,\log\!\frac{\rho+c}{\rho}+c\,\log\!\frac{\rho+c}{c}
  \;=\;(\rho+c)\,H_2\!\Bigl(\tfrac{\rho}{\rho+c}\Bigr),
  \label{eq:phi-star}
\end{equation}
where the last equality uses
$H_2(p)=-p\log p-(1-p)\log(1-p)$ with $p=\rho/(\rho+c)$.

\begin{lemma}[Pointwise surrogate gap on the disagreement set]
\label{lem:gap}
For every $x\in S$,
\begin{equation}
  \delta(x)\;\geq\;\frac{(\rho^\theta(x)-c(x))^2}{2\,(\rho^\theta(x)+c(x))}.
  \label{eq:gap-bound}
\end{equation}
\end{lemma}

\begin{proof}
Fix $x\in S$ and assume $\rho>c$; the case $\rho<c$ is symmetric and
the boundary case $\rho=c$ has zero right-hand side and is trivial.
Then $s^\star=\log(\rho/c)>0$ and, since $r_s(x)\neq r_{s^\star}(x)$,
$r_s(x)=0$, i.e.\ $s(x)\leq 0$. The function
$\Phi_{\mathrm{abs}}(\,\cdot\,;x)$ is strictly convex with unique
minimiser $s^\star>0$ (Proposition~\ref{prop:pointwise-min}); its
restriction to $(-\infty,0]$ is therefore strictly decreasing, and
attains its infimum at the right endpoint $s=0$. Hence
\begin{equation}
  \Phi_{\mathrm{abs}}(s(x);x)\;\geq\;\Phi_{\mathrm{abs}}(0;x),
  \qquad\text{so}\qquad
  \delta(x)\;\geq\;\Phi_{\mathrm{abs}}(0;x)-\Phi_{\mathrm{abs}}(s^\star;x).
  \label{eq:lb-zero}
\end{equation}
Combining \eqref{eq:phi-zero} and \eqref{eq:phi-star},
\begin{equation}
  \Phi_{\mathrm{abs}}(0;x)-\Phi_{\mathrm{abs}}(s^\star;x)
  \;=\;(\rho+c)\,\bigl[\log 2-H_2(p)\bigr]
  \;=\;(\rho+c)\,D_{\KL}\!\bigl(p\,\big\|\,\tfrac{1}{2}\bigr),
  \label{eq:phi-gap-kl}
\end{equation}
where $p:=\rho/(\rho+c)$ and the second equality uses the elementary
identity $D_{\KL}(p\,\|\,\tfrac{1}{2})=p\log(2p)+(1-p)\log(2(1-p))=\log
2-H_2(p)$. By Pinsker's inequality on
$\{0,1\}$~\citep{csiszar1967information} (in nats),
\[
  D_{\KL}\!\bigl(p\,\big\|\,\tfrac{1}{2}\bigr)\;\geq\;2\bigl(p-\tfrac{1}{2}\bigr)^2
  \;=\;\frac{(\rho-c)^2}{2(\rho+c)^2},
\]
where the last equality uses $p-\tfrac{1}{2}=(\rho-c)/(2(\rho+c))$.
Substituting into \eqref{eq:phi-gap-kl} and combining with
\eqref{eq:lb-zero} yields \eqref{eq:gap-bound}. The symmetric case
$\rho<c$ replaces ``decreasing on $(-\infty,0]$'' by ``increasing on
$[0,\infty)$'' but produces the same lower bound
$\Phi_{\mathrm{abs}}(0;x)-\Phi_{\mathrm{abs}}(s^\star;x)$.
\end{proof}

\begin{proof}[Proof of Theorem~\ref{thm:h-consistency}]
We first rewrite the abstention excess as an expectation over the
disagreement set $S$. By Proposition~\ref{prop:pointwise-min} and the
proof of Lemma~\ref{lem:bayes-rej},
$r_{s^\star}(x)=\ind\{\rho>c\}$ achieves the pointwise minimum
$\min\{\rho(x),c(x)\}$ of the proxy abstention loss
$\rho(x)\,\ind\{r(x)=0\}+c(x)\,\ind\{r(x)=1\}$. At $x\in S$ the rejector
$r_s$ chooses the opposite branch and pays $\max\{\rho(x),c(x)\}$
instead, an excess of $\max\{\rho,c\}-\min\{\rho,c\}=|\rho(x)-c(x)|$.
At $x\notin S$ the excess vanishes. Therefore
\begin{equation}
  L(r_s)-L(r_{s^\star})\;=\;\E_X\!\bigl[|\rho^\theta(X)-c(X)|\,\ind_{S}(X)\bigr].
  \label{eq:target-excess}
\end{equation}
Applying Cauchy--Schwarz to \eqref{eq:target-excess} with respect to the
measure $\Pr_X$, with factors
$f(x)=|\rho-c|\,\ind_{S}(x)$ and $g\equiv 1$,
\begin{equation}
  \bigl(L(r_s)-L(r_{s^\star})\bigr)^2
  \;\leq\;\E_X\!\bigl[(\rho^\theta(X)-c(X))^2\,\ind_{S}(X)\bigr].
  \label{eq:cs}
\end{equation}
By Lemma~\ref{lem:gap}, on $S$ we have
$(\rho-c)^2\leq 2(\rho+c)\,\delta(x)$, so
\begin{equation}
  \E_X\!\bigl[(\rho-c)^2\,\ind_{S}\bigr]
  \;\leq\;2\,\E_X\!\bigl[(\rho+c)\,\delta\,\ind_{S}\bigr]
  \;\leq\;2\,(\rho_{\max}+c_{\max})\,\E_X\!\bigl[\delta\,\ind_{S}\bigr]
  \;\leq\;2\,(\rho_{\max}+c_{\max})\,\E_X[\delta(X)],
  \label{eq:gap-agg}
\end{equation}
where the second inequality uses the essential-supremum bound
$\rho^\theta+c\leq\rho_{\max}+c_{\max}$ a.s.\ and the third uses
$\delta\geq 0$ everywhere on $\Omega$ (Proposition~\ref{prop:pointwise-min}).
By definition $\E_X[\delta(X)]=\Ecal_\Phi(s)-\Ecal_\Phi(s^\star)$;
substituting into~\eqref{eq:cs} via~\eqref{eq:gap-agg} and taking square
roots yields \eqref{eq:h-consistency-main}.
\end{proof}

\paragraph{On the constants and the rate.}
The square-root rate in \eqref{eq:h-consistency-main} is sharp for
strictly convex smooth surrogates and matches the calibration-function
rate of \citet{bartlett2006convexity}. The constant
$2(\rho_{\max}+c_{\max})$ is the natural one for our loss: it grows with
the heaviness of the right tail of $\rho^\theta$ but is partially
self-compensating, since the same outlier points produce large pointwise
$\delta(x)$ and therefore large $\Ecal_\Phi(s)-\Ecal_\Phi(s^\star)$ for
any $s$ that misclassifies them. Replacing the essential supremum by a
$\Pr_X$-quantile only weakens the bound by an additive
$\Pr_X(\rho^\theta+c>\tau)$ tail-mass term, which we do not pursue.

\subsection{Proof of Theorem~\ref{thm:hybrid-bound}}
\label{app:hybrid}

We use the discrete $\ell^2$ norm
$\|u\|_S^2:=\sum_{x\in S}\|u(x)\|^2$ over any subset $S\subseteq\Omega$,
which is additive over disjoint subsets: for disjoint $A,B\subseteq\Omega$,
$\|u\|_{A\cup B}^2=\|u\|_A^2+\|u\|_B^2$ and hence
$\|u\|_{A\cup B}\leq\|u\|_A+\|u\|_B$ via $\sqrt{a^2+b^2}\leq a+b$ for
$a,b\geq 0$. Recall the partition
$\Omega=\widehat{\Omega}_P\sqcup\widehat{\Omega}_C$ from
\eqref{eq:partition} with interface $\Gamma\subset\widehat{\Omega}_C$,
and the hybrid solution \eqref{eq:hybrid}.

Set $A:=\widehat{\Omega}_P\cup\Gamma$ and
$B:=\widehat{\Omega}_C\setminus\Gamma$, so $A\sqcup B=\Omega$. Applying
the disjoint-additivity inequality to
$\widehat{u}_{\mathrm{hyb}}-u^\star$,
\begin{equation}
  \bigl\|\widehat{u}_{\mathrm{hyb}}-u^\star\bigr\|_{\Omega}
  \;\leq\;
  \bigl\|\widehat{u}_{\mathrm{hyb}}-u^\star\bigr\|_{A}
  \;+\;
  \bigl\|\widehat{u}_{\mathrm{hyb}}-u^\star\bigr\|_{B}.
  \label{eq:split-norm}
\end{equation}

\paragraph{Term on $A$.}
By \eqref{eq:hybrid}, $\widehat{u}_{\mathrm{hyb}}(x)=h(x)$ for every
$x\in A$, so
\begin{equation}
  \bigl\|\widehat{u}_{\mathrm{hyb}}-u^\star\bigr\|_{A}
  \;=\;\bigl\|h-u^\star\bigr\|_{\widehat{\Omega}_P\cup\Gamma}.
  \label{eq:hyb-term-A}
\end{equation}

\paragraph{Term on $B$.}
On $B$, $\widehat{u}_{\mathrm{hyb}}=u_{\mathrm{slv}}$, where
$u_{\mathrm{slv}}$ is the fallback solution on $\widehat{\Omega}_C$ with
Dirichlet data $h\big|_\Gamma$. Introduce the auxiliary fallback
solution $\widetilde{u}_{\mathrm{slv}}$ on $\widehat{\Omega}_C$ driven
by the exact interface data $u^\star\big|_\Gamma$ and the same physical
boundary conditions on $\partial\widehat{\Omega}_C\cap\partial\Omega$.
By the triangle inequality,
\begin{equation}
  \bigl\|u_{\mathrm{slv}}-u^\star\bigr\|_{B}
  \;\leq\;
  \bigl\|u_{\mathrm{slv}}-\widetilde{u}_{\mathrm{slv}}\bigr\|_{B}
  \;+\;
  \bigl\|\widetilde{u}_{\mathrm{slv}}-u^\star\bigr\|_{B}.
  \label{eq:tri-inner}
\end{equation}
Both $u_{\mathrm{slv}}$ and $\widetilde{u}_{\mathrm{slv}}$ solve the
same discretised PDE on $\widehat{\Omega}_C$ with the same physical
boundary conditions and differ only in their Dirichlet interface data
($h\big|_\Gamma$ vs.\ $u^\star\big|_\Gamma$). Assumption~\ref{ass:stability}
applied to $g_1=h\big|_\Gamma$ and $g_2=u^\star\big|_\Gamma$ gives
\begin{equation}
  \bigl\|u_{\mathrm{slv}}-\widetilde{u}_{\mathrm{slv}}\bigr\|_{B}
  \;\leq\;
  C_s\,\bigl\|h-u^\star\bigr\|_{\Gamma}.
  \label{eq:stab-step}
\end{equation}
The second term in \eqref{eq:tri-inner} equals
$\epsilon_{\mathrm{disc}}$ by definition. Combining \eqref{eq:tri-inner}
and \eqref{eq:stab-step},
\begin{equation}
  \bigl\|\widehat{u}_{\mathrm{hyb}}-u^\star\bigr\|_{B}
  \;\leq\;
  C_s\,\bigl\|h-u^\star\bigr\|_{\Gamma}
  \;+\;\epsilon_{\mathrm{disc}}.
  \label{eq:hyb-term-B}
\end{equation}

\paragraph{Combining the two terms.}
Substituting \eqref{eq:hyb-term-A} and \eqref{eq:hyb-term-B} into
\eqref{eq:split-norm},
\begin{align*}
  \bigl\|\widehat{u}_{\mathrm{hyb}}-u^\star\bigr\|_{\Omega}
  &\;\leq\;
  \bigl\|h-u^\star\bigr\|_{\widehat{\Omega}_P\cup\Gamma}
  \;+\;C_s\,\bigl\|h-u^\star\bigr\|_{\Gamma}
  \;+\;\epsilon_{\mathrm{disc}}\\
  &\;\leq\;
  (1+C_s)\,\bigl\|h-u^\star\bigr\|_{\widehat{\Omega}_P\cup\Gamma}
  \;+\;\epsilon_{\mathrm{disc}},
\end{align*}
where the last inequality uses
$\|h-u^\star\|_{\Gamma}\leq\|h-u^\star\|_{\widehat{\Omega}_P\cup\Gamma}$
since $\Gamma\subseteq\widehat{\Omega}_P\cup\Gamma$. This is
\eqref{eq:hybrid-bound}. \qed

\section{Poisson experiment}
\label{app:poisson}

This appendix mirrors the Navier--Stokes study of
Section~\ref{sec:experiments} on a linear-elliptic Poisson problem.
The point of repeating the protocol on a different operator is
diagnostic: NSE is nonlinear, so the linearisation residual
$\delta_{\mathrm{lin}}^\theta$ in
Assumption~\ref{ass:linearisation} is non-trivial and the proxy
$\rho^\theta$ inherits that approximation slack. For the Poisson
operator $-\Delta$ the linearisation is exact -- $\Lcal_h=-\Delta$
regardless of $h$ -- and $e^\theta$ from \eqref{eq:ete} reconstructs
the pointwise PINN error to machine precision up to the
discretisation of $-\Delta^{-1}$. The Poisson experiment therefore
isolates the proxy from any nonlinear-coupling artefact, and the
abstention-loss-faithfulness conclusion of
Section~\ref{sec:exp-cylinder} should be even sharper here. We use
the same partitioning, the same surrogate \eqref{eq:surrogate}, the
same constant cost $c\equiv 1.1$, and the same TV-regularised CNN
class $\Hcal_r$ as in the cylinder experiment.

\subsection{Setup}

We solve the inhomogeneous Poisson equation $-\Delta u=f$ on the
unit square $\Omega=[0,1]^2$ with homogeneous Dirichlet boundary
data and a localised Gaussian source
\[
  f(x,y)\;=\;A\,\exp\!\Bigl(-\tfrac{(x-x_s)^2+(y-y_s)^2}{2\sigma^2}\Bigr),
\]
parameterised by source location $(x_s,y_s)$ and fixed amplitude and
width. The localisation gives the problem its diagnostic value: the
solution is sharply peaked near $(x_s,y_s)$ and smoothly decays
elsewhere, a regime in which PINNs reliably under-resolve the peak
neighbourhood and concentrate their pointwise error there. As in
Section~\ref{sec:exp-setup} we train a separate PINN
$h=\widehat{u}_\theta$ per source location, holding the predictor
fixed at evaluation. The fallback solver is linear FEM on a
structured triangulation of $\Omega$, which also produces the
reference $u^\star$. The error metric reduces to the scalar
$\mathrm{RMSE}(u,u^\star)=\sqrt{N^{-1}\sum_x(u(x)-u^\star(x))^2}$
because the pressure-gauge issues of the NSE setup do not arise
here.

The rejector is trained jointly on six source locations clustered
in the upper-right quadrant of $\Omega$ and evaluated on four
out-of-bracket source locations probing distinct generalisation
axes (Table~\ref{tab:poisson}). Each test source is paired with its
own dedicated PINN. Unlike the NSE setup, all four test sources are
deliberately \emph{outside} the training cluster: an above-cluster
source, an off-diagonal source, a far-field source in the
lower-left quadrant, and a grazing source just outside the
cluster. The Poisson experiment therefore tests the rejector's
behaviour under genuine extrapolation along the source-location
axis, which the in-bracket cylinder protocol does not.

\subsection{The abstention loss as the deployment-time metric}

The abstention-loss-faithfulness check that licensed reading
$\widehat{L}$ off the rejector at deployment in the cylinder case
applies a fortiori to Poisson. Because $\Lcal_h=-\Delta$ exactly,
the error-transport equation \eqref{eq:ete} recovers $h-u^\star$
without a linearisation residual, so $\rho^\theta$ tracks
$\|h-u^\star\|$ up to the discretisation error of the FFT inverse
of $-\Delta$. Sweeping $\tau$ on a held-out source location and
plotting $\widehat{L}$ and the true RMSE against coverage produces
the same U-shaped pair of curves as in Figure~\ref{fig:cyl-coverage},
with their minima co-locating at the same coverage. The qualitative
explanation is the one from Section~\ref{sec:exp-cylinder}: the
$\widehat{L}$ U comes from \eqref{eq:abst-empirical}, the RMSE U
comes from \eqref{eq:hybrid-bound}, and the alignment of the minima
is the empirical content of Theorem~\ref{thm:h-consistency} under
Assumption~\ref{ass:linearisation}. We therefore again read every
cross-configuration comparison off $\widehat{L}^\star$ and report
$\mathrm{RMSE}$ in Table~\ref{tab:poisson} as a redundancy check.

\subsection{Generalisation across source locations}

Table~\ref{tab:poisson} reports the abstention loss
$\widehat{L}^\star$ at the operating point $\tau^\star$, the
selected coverage, hybrid and PINN RMSE, and end-to-end wall-clock
time across all ten configurations, aggregated over three
independent rejector seeds (mean$\,\pm\,$std). The reading is the
one Section~\ref{sec:exp-generalisation} establishes for NSE, and
recurs here in a stronger form because the four test rows are
out-of-bracket rather than in-bracket. Neither failure mode appears:
$\widehat{L}^\star$ stays strictly inside the band bounded below by
$\E[\rho^\theta]$ (pure-accept) and above by $c=1.1$ (pure-abstain)
on every row, with values clustered in the $0.97\text{--}0.99$ band
across train and test alike, and coverages spread between $20\%$
and $50\%$. The four held-out source locations
(\texttt{ood\_pt\_above}, \texttt{ood\_pt\_offdiag},
\texttt{ood\_far}, \texttt{ood\_grazing}) sit in the same
$\widehat{L}^\star$ band as the training rows despite probing
genuinely distinct generalisation axes (above-cluster,
off-diagonal, far-field, grazing). A rejector that had memorised
the six training source locations would either spike
$\widehat{L}^\star$ on the held-out rows or collapse coverage to
the extremes; it does neither.

The RMSE columns confirm the same conclusion. Hybrid RMSE is
strictly below PINN RMSE on every row, and on each held-out source
the hybrid is $2\text{--}4\times$ below the PINN, indicating that
the rejector is not merely matching the PINN on unseen sources but
extracting genuine improvement by routing the peak neighbourhood to
FEM. End-to-end wall-clock is at near-parity with full FEM, which
is expected: the FEM solve on a $200\times200$ structured grid
takes $\approx 200$\,ms, so there is little runtime headroom for
the hybrid to recover -- the value of the hybrid here is in
accuracy and in the operational point that one can read
$\widehat{L}^\star$ at deployment without the FEM run.

\input{Sections/PoissonTable.tex}

The Poisson experiment thus reproduces the two structural claims
of Section~\ref{sec:experiments} on a linear elliptic operator and
under genuine source-location extrapolation: the abstention loss is
faithful (its operating point co-locates with the RMSE operating
point on a held-out source), and the rejector trained on six
clustered source locations transfers to four out-of-bracket sources
along distinct axes.
